\SourceLeaf{204}
\Article{7}
Wenn die Functionen, welche aus der Differentiation der willkürlichen Functionen \(f,f'\) entspringen, durch \(\varphi\) und \(\varphi'\) resp. bezeichnet werden, so dass \(\mathrm{d}\mathord{.}f\upsilon=\varphi\upsilon\mathord{.}\mathrm{d}\upsilon\), \(\mathrm{d}\mathord{.}f'\upsilon=\varphi'\upsilon\mathord{.}\mathrm{d}\upsilon\), so wird in Folge unsrer allgemeinen Auflösung
\[
\frac{\mathrm{d}P+i\,\mathrm{d}Q}{\mathrm{d}p+i\,\mathrm{d}q}=\varphi(p+iq),\qquad
\frac{\mathrm{d}P-i\,\mathrm{d}Q}{\mathrm{d}p-i\,\mathrm{d}q}=\varphi'(p-iq)
\]
\noindent also
\[
\frac{mmn}{N}=\varphi(p+iq)\mathord{.}\varphi'(p-iq)
\]
\noindent Das Vergrösserungsverhältniss bestimmt sich daher durch die Formel
\[
m=\BareRoot{\Bigl\{\frac{\mathrm{d}p^2+\mathrm{d}q^2}{\omega}\mathord{.}\frac{\Omega}{\mathrm{d}P^2+\mathrm{d}Q^2}\mathord{.}\varphi(p+iq)\mathord{.}\varphi'(p-iq)\Bigr\}}
\]

\Article{8}
Wir wollen nun noch unsre allgemeine Auflösung mit einigen Beispielen erläutern, wodurch sowohl die Art der Anwendung, als die Beschaffenheit einiger dabei noch in Betracht kommenden Umstände am besten ins Licht gesetzt werden wird.

Es seien zuvörderst beide Flächen Ebenen, wo wir
\[
\begin{aligned}
x&=t,&y&=u,&z&=0\\
X&=T,&Y&=U,&Z&=0
\end{aligned}
\]
\noindent werden setzen können. Die Differentialgleichung
\[
\omega=\mathrm{d}t^2+\mathrm{d}u^2=0
\]
\noindent gibt hier die beiden Integrale
\[
t+iu=\text{Const.},\quad t-iu=\text{Const.}
\]
\noindent und eben so sind die beiden Integrale der Gleichung \(\Omega=\mathrm{d}T^2+\mathrm{d}U^2=0\), folgende:
\[
T+iU=\text{Const.},\quad T-iU=\text{Const.}
\]
\noindent Die beiden allgemeinen Auflösungen der Aufgabe sind demnach:
\[
\begin{aligned}
\mathrm{I.}\quad&T+iU=f(t+iu),\quad T-iU=f'(t-iu)\\
\mathrm{II.}\quad&T+iU=f(t-iu),\quad T-iU=f'(t+iu).
\end{aligned}
\]
